\documentclass[11pt]{article}
\usepackage{mathblatt}
\newcommand{\bereich}[1]{\par\addvspace{14pt}\noindent{\large\bfseries #1}\par\nobreak\vspace{4pt}}
\newcommand{\bl}[1]{\par\noindent #1\par\vspace{3pt}}
\newcommand{\blrand}[1]{\par\noindent{\color{mbgrau}#1}\par\vspace{3pt}}
\newcommand{\blhier}[1]{\par\noindent\llap{$\blacktriangleright$\hspace{0.5em}}\textbf{#1}\par\vspace{3pt}}
\newcommand{\stufe}[1]{\par\noindent\hspace*{1.5em}{\small #1}\par\vspace{2pt}}
\newcommand{\pkt}{\textperiodcentered{}}
\begin{document}

\pfheftstil
\pflaufkopf[P10 \pkt{} Prüfung]{Dreiecke}{NKW}\pfseitevon
\pfheftkopf{Dreiecke}{}
\bereich{Rechtwinklige Dreiecke}
\bl{Länge mit Pythagoras}
\bl{Seite oder Winkel mit sin, cos, tan}
\bl{Seite im allgemeinen Dreieck (Sinussatz)}
\bereich{Alle Dreiecke und Vierecke}
\blrand{Vierecke: Symmetrie und Eigenschaften}
\blhier{Winkel ohne Rechnung bestimmen oder begründen}
\stufe{Winkelsumme}
\stufe{gleichschenkliges Dreieck}
\stufe{rechten Winkel begründen}
\bl{Symmetrie}

\par\addvspace{14pt}\noindent Weiter: Flächen\par
\end{document}